\section{Удлинение пути у Шоджи и в заметке о хроматичности}
\label{app:ch2path}

Этот разбор относится к $\Lambda$ формулы~\eqref{eq:beta-Lambda}. Период и знак $\delta_{\mathrm{eq}}$ выведены в разделе~\ref{sec:ch2}; здесь сравниваются два способа собрать то же удлинение.

\subsection{Среднее Шоджи}
\label{sec:ch2-shoji}

Шоджи~\cite{Shoji2005} начинает с разности пути за $n$ оборотов. Точная запись (4a) при обрезке выше второго порядка переходит в
\begin{equation}
\label{eq:beta-shoji-4b}
\Delta L_n
=\int_{s_0}^{s_0+n C_0}
\Biggl(hx+\frac{1}{2}x'^2+\frac{1}{2}y'^2\Biggr)\mathrm{d}s,
\end{equation}
где $h=1/\rho$. Среднее удлинение одного оборота есть предел (5), $\Delta C=\lim_{n\to\infty}\Delta L_n/n$. Это тот же предел, которым в разделе~\ref{sec:ch2-avg} отделена постоянная часть от колебания с бетатронной частотой.

Линейное $x_\beta=\sqrt{2J_x\beta_x}\sin\psi_x$ после суммирования по оборотам даёт формулы (10)--(12). Продольное смещение ограничено амплитудой $\sqrt{2J_x\mathcal{H}}$ и с числом оборотов не растёт. Постоянной части у линейного члена нет. Функция $\mathcal{H}=\gamma_x D^2+2\alpha_x DD'+\beta_x D'^2$ та же, что в разделе~\ref{sec:ch2-avg}.

Вклад наклонов --- формулы (16) и (17). Производная записана как $y'=\sqrt{2J_y\gamma_y}\sin\psi_y$. Такая одночленная запись есть фазовый сдвиг точной производной: амплитуда $\sqrt{2J\gamma}$ имеет тот же средний квадрат $J\gamma$, что и~\eqref{eq:beta-xp2}, и для неё не требуется $\alpha=0$. После отбрасывания осциллирующей части
\begin{equation}
\label{eq:beta-shoji-slope}
\lim_{n\to\infty}\frac{1}{n}\int\frac{1}{2}y'^2\,\mathrm{d}s
=\frac{J_y}{2}\int_0^{C_0}\gamma_y\,\mathrm{d}s,
\end{equation}
и то же для $x$ с $J_x$ и $\gamma_x$. Деление на $C_s$ даёт угловую часть~\eqref{eq:beta-Lambda}.

Нелинейный сдвиг $x$ собирается из шести квадратичных членов правой части уравнения (2a): секступольного $m$, кривизны $h^3$ и краевого $h'$. Для одного тонкого толчка $\Delta\theta=\Delta\kappa\, y^2$ среднее по оборотам есть формула (24),
\begin{equation}
\label{eq:beta-shoji-24}
\Delta C=\Delta\kappa\, J_y\beta_y D.
\end{equation}
Множитель $1/2$ в эту строку не входит, потому что он уже сидит в коэффициенте $\Delta\kappa$, когда толчок берётся из уравнения движения: там $\Delta\kappa=(1/2)(h''-hk-m)\,\mathrm{d}s$. Средний квадрат $\langle y^2\rangle=J_y\beta_y$ тогда даёт в длине $(J_y/2)\int m\beta_y D\,\mathrm{d}s$ с тем знаком, какой у $m$ в уравнении. Остальные пять членов считаются той же подстановкой; их список вынесен в приложение~A статьи. Сумма с угловыми членами (16) и (17) есть полная формула (26).

Хроматичности получаются из тех же уравнений (2) заменой $x\to x+D\delta$. Формулы (29) интегрированием по частям, тождества которого собраны в приложении~B, переписываются как (30). Сравнение (30) с (26) сокращает подынтегральные выражения до
\begin{equation}
\label{eq:beta-shoji-31}
\Delta C=-2\pi(J_x\xi_x+J_y\xi_y).
\end{equation}
Деление на $C_s$ есть~\eqref{eq:beta-xi}. Равенство удерживает секступоль, кривизну и краевое поле: они стоят и в пути, и в хроматичности. Короткий секступоль для этого вывода не требуется. Шоджи отмечает, что формула такого вида без вывода была у Эмери, а Форест связал её с симплектичностью. Выше она получена из интеграла пути, и на это замечание вывод не опирается.

Для электронов Шоджи закрывает удлинение сдвигом импульса по формуле (32),
\begin{equation}
\label{eq:beta-shoji-32}
\delta=-\frac{1}{\alpha}\frac{\Delta C}{C_0}.
\end{equation}
Делитель $\alpha$ --- коэффициент уплотнения, не $\eta$. Для электронов на $1$~ГэВ величина $1/\gamma^2$ мала рядом с обычным $\alpha$ порядка $10^{-2}$ и меньше; в квазиизохронной строке таблицы~I той работы $\alpha=2\times 10^{-6}$ уже сравним с $1/\gamma^2$. Для дейтронов магнитного кольца знаменатель равновесия есть $\eta_0=-0{,}534211$, тогда как $\alpha_0=0{,}230$. Подстановка $\alpha_0$ в~\eqref{eq:beta-deq} меняет и величину, и знак $\delta_{\mathrm{eq}}$. Формула (32) этот знаменатель не заменяет.

\subsection{Заметка о хроматичности}
\label{sec:ch2-note}

Заметка~\cite{SenichevNote} идёт от геометрического интеграла (7),
\begin{equation}
\label{eq:beta-note-7}
\Delta C
=\int\Biggl(\frac{x}{\rho}+\frac{x'^2+y'^2}{2}\Biggr)\mathrm{d}s.
\end{equation}
Это та же обрезка, что~\eqref{eq:beta-dl-2} и~\eqref{eq:beta-shoji-4b}. Уравнение движения (1) заметки есть $x''+K_1 x=(1/\rho)\delta+F^x$ и $y''-K_1 y=F^y$, с секступольными силами $F^x=(K_2/2)(x^2-y^2)$ и $F^y=-K_2 xy$.

Формула (8) ещё содержит $\alpha$:
\begin{equation}
\label{eq:beta-note-8}
x'_\beta
=-\sqrt{\frac{2J_x}{\beta_x}}
\bigl(\alpha_x\cos\psi_x+\sin\psi_x\bigr).
\end{equation}
Следующая фраза считает среднее $\beta'$ по кольцу равным нулю и заключает, что $\alpha=0$. Из периодичности $\beta$ следует только $\int\beta'\,\mathrm{d}s=0$. Локально $\alpha=-\beta'/2$ нулём не является. Формула (9) после этого шага пишет
\begin{equation}
\label{eq:beta-note-9}
x'_\beta
=-\sqrt{\frac{2J_x}{\beta_x}}\sin\psi_x
=-\sqrt{2J_x\gamma_x}\,\sin\psi_x.
\end{equation}
Равенство двух амплитуд требует $\gamma=1/\beta$, то есть $\alpha=0$ в той же точке. Оно ложно, пока $\alpha\neq 0$. Подстановка первой амплитуды в~\eqref{eq:beta-note-7} дала бы $(J/2)\int\mathrm{d}s/\beta$ и потеряла бы $(J/2)\int(\alpha^2/\beta)\,\mathrm{d}s$. В интеграл заметка подставляет вторую амплитуду. Средний квадрат $\langle x'^2\rangle=J\gamma$ у неё верный, и формула (11) заметки есть
\begin{equation}
\label{eq:beta-note-11}
\Delta C_{x',y'}
=\frac{J_x}{2}\int\gamma_x\,\mathrm{d}s
+\frac{J_y}{2}\int\gamma_y\,\mathrm{d}s.
\end{equation}
Это~\eqref{eq:beta-angle}. Тождество~\eqref{eq:beta-xp2} удерживает $\alpha$ внутри квадрата и к шагу $\alpha=0$ не обращается; конечная формула (11) с ним совпадает. Линейный член $x_\beta/\rho$ в среднем по кольцу равен нулю, как и у Шоджи.

Дисперсионная часть записана как $\delta(\alpha_0+\alpha_1\delta+\cdots)$, формулы (13) и (14). Коэффициент $\alpha$ там есть интеграл $\int D/\rho\,\mathrm{d}s$ и имеет размерность длины. Безразмерные $\alpha_0$ и $\alpha_1$ формулы~\eqref{eq:mag-C} получаются делением на длину орбиты. Постоянный толчок даёт формулу (15), $\Delta C=\sum F_{\mathrm{kick}} L D$. В секступоле квадрат отклонения заменяется средним по бетатронной фазе. Из $\langle x^2\rangle=J_x\beta_x$ и $\langle y^2\rangle=J_y\beta_y$ сила $(K_2/2)(x^2-y^2)$ даёт
\begin{equation}
\label{eq:beta-note-kick}
\langle F^x\rangle L
=\frac{1}{2}K_2 L(J_x\beta_x-J_y\beta_y).
\end{equation}
Множитель $1/2$ здесь стоит в определении силы, а второй $1/2$ возникает из среднего $\cos^2\psi$. Произведение совпадает с упаковкой Шоджи, где та же $1/2$ спрятана в $\Delta\kappa$ формулы~\eqref{eq:beta-shoji-24}, а $\langle y^2\rangle=J\beta$ стоит без дополнительной половины. Сумма углового члена, дисперсии и секступолей есть формула (20) заметки,
\begin{equation}
\label{eq:beta-note-20}
\begin{aligned}
\Delta C
&=\frac{J_x}{2}\int(\gamma_x+K_2\beta_x D)\,\mathrm{d}s
+\frac{J_y}{2}\int(\gamma_y-K_2\beta_y D)\,\mathrm{d}s
\\
&\quad+\delta(\alpha_0+\alpha_1\delta+\cdots).
\end{aligned}
\end{equation}
Кривизны и краевого поля в этой сумме нет. Короткий секступоль --- условие, при котором нелинейный сдвиг заменён средним толчком~\eqref{eq:beta-note-kick}.

Хроматичность берётся в виде (21),
\begin{equation}
\label{eq:beta-note-21}
\begin{aligned}
\xi_x
&=-\frac{1}{4\pi}\int(K_1\beta_x+K_2\beta_x D)\,\mathrm{d}s,
\\
\xi_y
&=-\frac{1}{4\pi}\int(K_1\beta_y-K_2\beta_y D)\,\mathrm{d}s.
\end{aligned}
\end{equation}
Горизонтальная строка следует из уравнения $x''+K_x x=0$ с $K_x=K_1-K_2 D\delta$ и из правила $\mathrm{d}\nu=(1/(4\pi))\int\beta\,\mathrm{d}K$. Масштаб градиента $K_1/(1+\delta)$ даёт $-K_1$, секступольный толчок даёт $-K_2 D$. Дипольного $1/\rho^2$ в строке нет: естественная горизонтальная хроматичность содержит $K_x=K_1+1/\rho^2$. Вертикальное уравнение $y''-K_1 y=-K_2 xy$ имеет фокусирующий коэффициент $K_y=-K_1$. Его производная по $\delta$ равна $K_1+K_2 D$, поэтому
\begin{equation}
\label{eq:beta-xi-vert}
\xi_y
=\frac{1}{4\pi}\int\beta_y(K_1+K_2 D)\,\mathrm{d}s.
\end{equation}
Знак $K_2$ в~\eqref{eq:beta-note-21} с этой формулой совпадает. Знак $K_1$ противоположен: в~\eqref{eq:beta-note-21} перед $K_1\beta_y$ стоит тот же минус, что и в горизонтальной строке.

Тождество Твисса $\alpha'=K\beta-\gamma$ на периодическом кольце даёт $\int K\beta\,\mathrm{d}s=\int\gamma\,\mathrm{d}s$, где $K$ --- коэффициент при $z$ в $z''+Kz=0$. Для горизонтали $K_x=K_1+1/\rho^2$, для вертикали $K_y=-K_1$. Формула (22) заметки пишет
\begin{equation}
\label{eq:beta-note-22}
\int K_1\beta_{x,y}\,\mathrm{d}s=\int\gamma_{x,y}\,\mathrm{d}s
\end{equation}
сразу для обеих плоскостей. По горизонтали выпало $\int\beta_x/\rho^2\,\mathrm{d}s$. По вертикали знак $K_1$ противоположен $K_y$, и верное равенство есть $\int(-K_1)\beta_y\,\mathrm{d}s=\int\gamma_y\,\mathrm{d}s$.

Если подставить~\eqref{eq:beta-note-22} в~\eqref{eq:beta-note-21}, квадрупольный кусок становится $\xi=-(1/(4\pi))\int\gamma\,\mathrm{d}s$ в обеих плоскостях. Для вертикали это получается из двух неверных знаков сразу: минус перед $K_1$ в (21) и плюс в $\int K_1\beta_y=\int\gamma_y$. Вместе они дают то $\xi_y$, которое переводит~\eqref{eq:beta-note-20} в
\begin{equation}
\label{eq:beta-note-24}
\Delta C
=-2\pi(J_x\xi_x+J_y\xi_y)
+\delta(\alpha_0+\alpha_1\delta+\cdots).
\end{equation}
Секступольные знаки в (21) с толчком~\eqref{eq:beta-note-kick} согласованы и в этой компенсации не участвуют. Бетатронная часть~\eqref{eq:beta-note-24} поэтому совпадает с путём~\eqref{eq:beta-note-20}. С формулой Шоджи~\eqref{eq:beta-shoji-31} она совпадает на квадрате угла и на коротком секступоле. Кривизна, краевое поле и горизонтальное $1/\rho^2$ в (24) не входят: первые два не удерживались в пути, третье отсутствует в обеих строках (21) и (22).

Вывод заметки состоит в том, что при коротких секступолях нулевая хроматичность обнуляет бетатронное удлинение. В пределах~\eqref{eq:beta-note-20} это верно. Полное $\Lambda$ обнуляется вместе с $\xi_x$ и $\xi_y$ и в модели Шоджи, где кривизна и край уже включены в обе части (31). Условие $\eta_1=0$ к этой заметке не относится: её $\alpha$ есть размерный интеграл длины, а равновесие периода читается из~\eqref{eq:beta-sync}.
